\documentclass[letterpaper]{article} % DO NOT CHANGE THIS
\usepackage[submission]{aaai2027} % DO NOT CHANGE THIS
\usepackage[hyphens]{url} % DO NOT CHANGE THIS
\usepackage{graphicx} % DO NOT CHANGE THIS
\urlstyle{rm} % DO NOT CHANGE THIS
\def\UrlFont{\rm} % DO NOT CHANGE THIS
\usepackage{natbib} % DO NOT CHANGE THIS AND DO NOT ADD ANY OPTIONS TO IT
\usepackage{caption} % DO NOT CHANGE THIS AND DO NOT ADD ANY OPTIONS TO IT
\frenchspacing % DO NOT CHANGE THIS

\usepackage{algorithm}
\usepackage{algorithmic}

\usepackage{newfloat}
\usepackage{listings}
\DeclareCaptionStyle{ruled}{labelfont=normalfont,labelsep=colon,strut=off}
\lstset{
    basicstyle={\footnotesize\ttfamily},
    numbers=left,
    numberstyle=\footnotesize,
    xleftmargin=2em,
    aboveskip=0pt,
    belowskip=0pt,
    showstringspaces=false,
    tabsize=2,
    breaklines=true
}
\floatstyle{ruled}
\newfloat{listing}{tb}{lst}{}
\floatname{listing}{Listing}

\usepackage{booktabs}

\pdfinfo{
/TemplateVersion (2027.1)
}

\setcounter{secnumdepth}{0}

\title{Trabalho prático 1: relatório}
\author{
    João Gabriel Vaz
}
\affiliations{
    [Instituição]\\
    [Curso ou Disciplina]\\
    [Matrícula]\\
    [Endereço de E-mail]
}

\begin{document}

\maketitle

\section{Introdução e Objetivo}

O Ligue-4 é um jogo determinístico, de soma zero e informação perfeita, cuja árvore de estados torna inviável a exploração exaustiva durante cada jogada. Este trabalho desenvolve um agente capaz de escolher movimentos por busca competitiva sob restrições de profundidade e tempo. O objetivo é estudar a evolução do agente desde o Minimax com avaliação heurística até a poda Alfa--Beta e o \textit{Iterative Deepening}, comparando taxa de vitória, tempo de decisão, estados visitados e profundidade alcançada. Por fim, uma versão competitiva incorpora técnicas adicionais para aproveitar o limite de três segundos por jogada sem comprometer a validade da resposta.

\section{Metodologia}

\paragraph{Representação e Minimax.}
Um estado é representado por uma matriz $6 \times 7$, e seus sucessores são produzidos pela inserção de uma peça, respeitando a gravidade, em cada coluna não cheia. A perspectiva de maximização permanece fixa no jogador para o qual \texttt{choose\_move} foi chamado. Estados de vitória e derrota recebem, respectivamente, $10^6+d_r$ e $-(10^6+d_r)$, onde $d_r$ é a profundidade restante. Dessa forma, o agente prefere vitórias mais rápidas e, quando a derrota é inevitável, procura adiá-la. Empates recebem valor zero. Ao atingir o limite de profundidade em um estado não terminal, aplica-se a função heurística. As jogadas seguem inicialmente a ordem $(3,2,4,1,5,0,6)$, priorizando o centro e estabelecendo desempates determinísticos.

\paragraph{Função heurística.}
A avaliação é calculada da perspectiva do jogador maximizador por $h(s)=6(C_{\mathrm{max}}-C_{\mathrm{min}})+\sum_{w\in W}g(w)$, em que $C_{\mathrm{max}}$ e $C_{\mathrm{min}}$ são as quantidades de peças dos jogadores na coluna central e $W$ contém as 69 janelas horizontais, verticais e diagonais de quatro casas. Janelas ocupadas pelos dois jogadores recebem zero, pois já não podem formar uma vitória. Para janelas exclusivas do agente, uma peça vale $+1$ e duas valem $+10$; para o adversário, os valores correspondentes são $-1$ e $-12$. Três peças com uma casa vazia valem $+500$ quando essa casa já é jogável e $+40$ caso contrário. A configuração adversária equivalente vale $-650$ ou $-50$. Uma casa é considerada jogável quando está na última linha ou possui apoio imediatamente abaixo. A assimetria das penalidades torna o agente mais conservador diante de ameaças adversárias imediatas.

\paragraph{Poda e controle de tempo.}
Para uma mesma profundidade e sem interrupção temporal, a poda Alfa--Beta preserva o valor do Minimax, mas elimina ramos incapazes de modificar a decisão. O \textit{Iterative Deepening} executa buscas Alfa--Beta completas nas profundidades $1,2,\ldots$, até o limite configurado ou até o prazo. O relógio é monotônico e reserva uma margem de segurança de $20\%$ do orçamento, limitada a $300$\,ms. Somente uma iteração integralmente concluída substitui a melhor jogada conhecida; movimento, valor e dicas da iteração interrompida não são promovidos. Na Alfa--Beta de profundidade fixa, somente filhos da raiz completamente avaliados podem determinar a resposta. Se nenhum resultado utilizável estiver disponível, aplica-se um \textit{fallback} que procura, nesta ordem, vitória imediata, a primeira ameaça adversária imediata segundo a ordem central e a coluna legal mais central. O contador de estados inclui raiz, nós internos e folhas; no aprofundamento iterativo, acumula todas as iterações, inclusive a parte explorada da tentativa interrompida.

\paragraph{Versão competitiva.}
A versão final combina \textit{Iterative Deepening}, Alfa--Beta e uma tabela de transposição renovada a cada profundidade. A chave reúne o tabuleiro, o jogador corrente, o jogador maximizador e a profundidade restante. Cada entrada armazena valor, tipo de limite e, quando aplicável, o melhor movimento; os tipos são \textsc{Exact}, \textsc{Lower} e \textsc{Upper}, e a tabela admite até 200 mil entradas. Dicas de movimento só são promovidas após uma iteração completa. A ordenação prioriza vitórias e bloqueios imediatos, o movimento da tabela, a dica anterior e, então, a ordem central. Na raiz, vitórias diretas e bloqueios forçados únicos são resolvidos sem busca. Os filhos são gerados uma vez e os empates preservam a prioridade central canônica, aumentando a profundidade efetiva sem alterar as versões básicas dos experimentos.

\section{Experimentos e Resultados}

\paragraph{Protocolo.}
Foram executadas quinze partidas, sempre com o primeiro algoritmo do confronto como jogador~1. Nos Experimentos~1 e~2, utilizaram-se profundidades $D=2,\ldots,5$ e limite de $30\,000$\,ms. No Experimento~3, o limite máximo foi $D=42$, com orçamentos de 1 e 2\,s. Por fim, foram disputadas cinco partidas humanas contra a versão competitiva, com $T=3$\,s e limites $D=7,6,8,9,10$, respectivamente. Cada configuração automática foi executada uma vez; portanto, as taxas de vitória são apenas descritivas. O tempo reportado é o tempo interno de busca, sem a sobrecarga do servidor, e $N$ denota estados visitados por decisão. Nenhuma partida utilizou o \textit{fallback} externo do servidor.

\begin{table}[t]
\centering
\caption{Resultados com profundidade fixa. Em cada linha, $A/B$ segue a ordem do confronto; $n$ é o número de decisões e $C_A$ são os cortes do primeiro algoritmo.}
\label{tab:fixed-depth}
\scriptsize
\resizebox{\columnwidth}{!}{
\begin{tabular}{lrrrrrrr}
\toprule
$A$--$B$ & $D$ & Res. & L & $n_A/n_B$ &
$t_A/t_B$ (ms) & $N_A/N_B$ & $C_A$ \\
\midrule
M--R  & 2 & M  & 7  & 4/3   & 4,00/0,01     & 53,5/0          & -- \\
M--R  & 3 & M  & 7  & 4/3   & 26,49/0,01    & 386,0/0         & -- \\
M--R  & 4 & M  & 7  & 4/3   & 176,32/0,01   & 2454,5/0        & -- \\
M--R  & 5 & M  & 7  & 4/3   & 1114,41/0,01  & 17031,3/0       & -- \\
AB--M & 2 & AB & 39 & 20/19 & 1,47/2,34     & 19,5/32,9       & 2,80 \\
AB--M & 3 & E  & 42 & 21/21 & 5,96/13,61    & 90,0/196,7      & 9,86 \\
AB--M & 4 & AB & 29 & 15/14 & 15,40/107,07  & 260,7/1580,2    & 57,13 \\
AB--M & 5 & M  & 38 & 19/19 & 127,02/738,79 & 2190,2/11112,8  & 339,32 \\
\bottomrule
\end{tabular}}
\end{table}

\paragraph{Experimentos 1 e 2: profundidade fixa.}
O Minimax venceu o agente aleatório nas quatro profundidades, totalizando $4/4$ vitórias. Entretanto, entre $D=2$ e $D=5$, sua média de estados cresceu aproximadamente 318 vezes e o tempo, 279 vezes. No segundo confronto, a Alfa--Beta obteve duas vitórias, um empate e uma derrota, correspondendo a taxas observadas de $50\%$, $25\%$ e $25\%$, respectivamente. A poda reduziu os estados visitados entre $40,8\%$ e $83,5\%$ e o tempo entre $37,3\%$ e $85,6\%$. Essas reduções comparam médias de turnos distintos, pois os algoritmos ocuparam lados opostos, mas demonstram claramente a diferença de custo. A vitória do Minimax em $D=5$ não contradiz essa eficiência: a poda preserva o valor da busca completa, mas não garante o resultado de uma partida específica.

\begin{table}[t]
\centering
\caption{\textit{Iterative Deepening} (ID) contra Alfa--Beta direta (AB), com $D=42$. $d_{\mathrm{ID}}$ é a profundidade completa média e $F_{\mathrm{AB}}$ conta \textit{fallbacks} internos.}
\label{tab:timed}
\scriptsize
\resizebox{\columnwidth}{!}{
\begin{tabular}{rrrrrrrr}
\toprule
$T$ (s) & Res. & L & $n_{\mathrm{ID}}/n_{\mathrm{AB}}$ &
$t_{\mathrm{ID}}/t_{\mathrm{AB}}$ (ms) &
$d_{\mathrm{ID}}$ & $N_{\mathrm{ID}}/N_{\mathrm{AB}}$ &
$F_{\mathrm{AB}}$ \\
\midrule
1 & ID & 11 & 6/5 & 800,07/800,05   & 6,33 & 14088/66478   & 5/5 \\
2 & ID & 11 & 6/5 & 1700,06/1700,05 & 7,00 & 31081/141844  & 5/5 \\
\bottomrule
\end{tabular}}
\end{table}

\paragraph{Experimento 3: busca temporizada.}
O ID venceu ambas as partidas, isto é, $2/2$ vitórias observadas. Com 2\,s, sua profundidade completa média aumentou de $6,33$ para $7,00$. Todas as decisões atingiram o prazo interno, mas o ID sempre reutilizou uma iteração concluída. A Alfa--Beta direta não terminou nenhum filho da raiz em qualquer decisão e, por isso, acionou em $10/10$ casos o \textit{fallback} tático interno. Assim, o confronto evidencia principalmente a propriedade \textit{anytime} do ID, e não uma comparação entre duas buscas completas. Os tempos de aproximadamente 800 e 1700\,ms decorrem da margem reservada para que a resposta alcance o servidor.

\begin{table}[t]
\centering
\caption{Versão competitiva contra humano. Todas as partidas foram vencidas pela IA. $\bar t$ e $\bar N$ incluem todas as decisões; $\bar d_s$ considera somente buscas, excluindo o atalho final.}
\label{tab:human}
\scriptsize
\resizebox{\columnwidth}{!}{
\begin{tabular}{rrrrrrrr}
\toprule
Jogo & $D$ & L & $n_{\mathrm{IA}}$ &
$\bar t$ (ms) & $\bar N$ & $\bar d_s$ & TO \\
\midrule
1 & 7  & 9  & 5 & 341,79  & 5823  & 7,00 & 0 \\
2 & 6  & 17 & 9 & 176,97  & 2835  & 6,00 & 0 \\
3 & 8  & 9  & 5 & 751,00  & 13303 & 8,00 & 0 \\
4 & 9  & 15 & 8 & 1792,68 & 34607 & 8,71 & 2 \\
5 & 10 & 7  & 4 & 2027,02 & 35352 & 8,33 & 3 \\
\bottomrule
\end{tabular}}
\end{table}

\paragraph{Experimento 4: versão competitiva.}
A IA venceu as cinco partidas humanas, totalizando $5/5$ vitórias e 57 lances. Cada partida terminou com o reconhecimento de uma vitória imediata, sem \textit{fallback} interno. Os limites foram integralmente alcançados entre $D=6$ e $D=8$; em $D=9$ ocorreram dois \textit{timeouts} internos e, em $D=10$, três. A profundidade completa média cair de $8,71$ para $8,33$ entre esses dois jogos mostra a saturação imposta pelo orçamento, não uma regressão do algoritmo. Como posições, jogadas humanas e limites variaram, os comprimentos das partidas não constituem uma comparação controlada entre profundidades.

\section{Discussão}

O crescimento observado no Minimax confirma o custo exponencial da exploração sem poda. A Alfa--Beta reduziu consistentemente tempo e estados, com maior efeito nas profundidades mais altas, sem modificar a decisão que seria obtida pelo Minimax sobre o mesmo estado e profundidade. Entretanto, como os dois algoritmos jogaram lados opostos, as razões da Tabela~\ref{tab:fixed-depth} devem ser interpretadas como indicadores experimentais, não como uma medição pareada.

Sob prazo rígido, o \textit{Iterative Deepening} mostrou-se mais robusto: uma interrupção conserva o resultado da última profundidade concluída. A busca Alfa--Beta direta em $D=42$, por outro lado, consumiu todo o intervalo sem concluir um filho da raiz, fazendo com que sua decisão dependesse do \textit{fallback}. O teste é válido para avaliar a política de interrupção definida, mas favorece a análise da propriedade \textit{anytime}, e não da força relativa dos algoritmos em profundidades concluídas. A margem interna evitou substituições pelo servidor, ao custo de reservar 200--300\,ms do orçamento.

A versão competitiva combinou essa robustez com ordenação tática, tabela de transposição e atalhos imediatos. Todas as 26 buscas substanciais das partidas humanas registraram acertos na tabela, confirmando seu uso; contudo, sem um experimento de ablação, não é possível isolar seu ganho. Os resultados também são limitados pelo pequeno número de partidas: houve uma execução por configuração automática, o algoritmo avaliado sempre iniciou, o agente aleatório não teve semente controlada e as partidas humanas usaram profundidades e posições diferentes. Assim, as taxas de vitória descrevem apenas esta amostra e não constituem estimativas estatísticas gerais.

\section{Conclusão}

Os experimentos indicam que a poda Alfa--Beta é substancialmente mais eficiente que o Minimax em profundidade fixa e que o \textit{Iterative Deepening} é mais adequado a limites temporais, pois sempre preserva uma decisão de busca completa. A versão competitiva respondeu sem \textit{fallback} externo e venceu as cinco partidas humanas observadas, embora $D=9$ e $D=10$ já tenham revelado saturação temporal.

Como continuidade, recomenda-se repetir os confrontos com alternância de quem inicia, avaliar os algoritmos sobre os mesmos estados e realizar ablações da heurística, ordenação e tabela de transposição. Representações por \textit{bitboards}, avaliação incremental e ajuste dos pesos por autojogo também podem ampliar a profundidade atingida.

\end{document}
